% TOP_PROBLEM: {"problemId": "problem.kothe-conjecture-on-nil-one-sided-ideals", "canonicalTitle": "Köthe Conjecture on Nil One-Sided Ideals", "releaseRank": 467, "primaryDomain": "algebra_representation_category", "primaryDomainLabel": "Algebra, representation & category theory", "displayStatus": "open", "releaseStatus": "open"}
% !TeX program = lualatex
% ATTEMPT_STATUS: unresolved
\documentclass[11pt]{article}
\usepackage[margin=25mm]{geometry}
\usepackage{fontspec}
\setmainfont{Latin Modern Roman}
\usepackage{amsmath,amssymb,mathtools}
\usepackage{unicode-math}
\setmathfont{Latin Modern Math}
\usepackage{xurl,hyperref}
\hypersetup{colorlinks=true,urlcolor=blue}
\setcounter{secnumdepth}{0}
\setlength{\parindent}{0pt}
\setlength{\parskip}{5pt}
\setlength{\emergencystretch}{3em}
\begin{document}
\section*{\texorpdfstring{Köthe Conjecture on Nil One-Sided Ideals}{Köthe Conjecture on Nil One-Sided Ideals}}\label{467}
\subsection{Definitions and mathematical statement}
An associative ring or ideal is nil if each of its elements is nilpotent, with a nilpotence exponent that may depend on the element. This is weaker than being nilpotent, which requires one uniform bound on lengths of arbitrary products. Rings in the nil-ring formulations need not have an identity.

K\"othe's assertion is that a ring with no nonzero nil two-sided ideal also has no nonzero nil one-sided ideal. The equivalent formulations recorded in the source include: the sum of two nil right ideals is nil; and
\[
 R\text{ nil}\quad\Longrightarrow\quad M_2(R)\text{ nil}.
\]
The conclusion in the matrix version concerns every matrix individually, not a uniform common nilpotence exponent. Commutativity and finite dimensionality are not hypotheses.
\par\smallskip\textit{Formulation and concept references:} \href{https://www.mathunion.org/fileadmin/ICM/Proceedings/ICM2006.2/ICM2006.2.ocr.pdf}{[S1]}.

\subsection{Short English statement}
Can nil behavior in a one-sided ideal fail to produce a nonzero nil two-sided ideal? Equivalently, do two-by-two matrices over nil rings remain nil?
\subsection{Research attempt}
\paragraph{A current-status qualification comes first.}
This attempt was prepared on 27 September 2026. The inherited catalogue
metadata says ``open'', but the source audit found a directly relevant
new primary manuscript: Adamczewski--B\"ohmler--Marczinzik [E1],
arXiv:2609.07996v1, submitted 7 September 2026, manuscript dated
5 September 2026. Its Theorem 1.1 asserts that for every countable field
$F$ there is a nil algebra $N=F\langle b,c\rangle_+$ for which
\[
 W=\begin{pmatrix}b^2+c&bc\\ b&c\end{pmatrix}
 \quad\text{is not nilpotent in }M_2(N).
\]
Thus the external claim addresses the exact selected matrix assertion,
not merely a polynomial-ring analogue. The visible arXiv record showed
only version 1 and no withdrawal notice or qualification of that theorem;
a targeted correction search found none. This is not a certification of
the complete new proof or a prediction of its eventual reception.

The \texttt{unresolved} status of this document refers to the present
attempt's own result. It must not be read as an assertion that the
literature frontier is unchanged. The inherited definition and source
entry are retained, while the new claim is explicitly recorded. The
original large ICM PDF [S1] was not retrievable in this audit; [E2] was
checked for the classical nil-ideal formulation. The work below gives
independently justified reductions and subcases rather than presenting
the external construction as our own proof.

\paragraph{Keep three quantifiers separate.}
An element $a$ is nilpotent when there exists $e(a)$ with $a^{e(a)}=0$.
A ring is nil when this holds for every element. A ring is nilpotent when
one $r$ annihilates every product of $r$ possibly different elements.
Local nilpotence asks that every subring generated by finitely many elements
be nilpotent, with a bound allowed to depend on that finite set.

The matrix question starts with nilness only. A matrix has finitely many
entries, but its powers contain products of those entries in many orders.
Replacing nilness by local nilpotence would supply exactly the missing
uniform control on these mixed products. This replacement is not licensed
by the statement, and a purported proof making it would already have
assumed a stronger property of the input ring.

All the subring constructions below omit a new identity. When convenient
we embed a ring in a unitization for module arguments; a nonzero nil ring
itself cannot have an identity, since that identity would not be nilpotent.

\paragraph{Lemma 467.1: the finite-entry reduction and a complete subcase.}
Let $A\in M_d(R)$, and let $B\subseteq R$ be the subring generated by
its $d^2$ entries. If $B^r=0$, then $A^r=0$. Consequently all matrices
over a locally nilpotent ring are nilpotent individually.

\emph{Proof.} The $(i,j)$ entry of $A^r$ is
\[
 \sum_{i_1,\ldots,i_{r-1}=1}^d
 A_{ii_1}A_{i_1i_2}\cdots A_{i_{r-1}j}.
\]
Every summand is a product of $r$ elements of $B$ and hence vanishes.
For a locally nilpotent ring, the particular finitely generated subring
$B$ has such an exponent. The exponent may vary from matrix to matrix.
\hfill$\square$

The same formula gives a useful countability reduction. If a matrix over
some nil ring is not nilpotent, the subring generated by its finitely many
entries is still nil, is countable, and contains the same nonnil matrix.
Indeed its elements are integer linear combinations of finite nonempty
words in a finite alphabet, a countable set. This does not reduce an
arbitrary ring to an algebra over a field: its characteristic and additive
torsion must still be respected. It does reduce the number of generators
needed to study a proposed matrix witness.

\paragraph{Lemma 467.2: the commutative case, with an explicit exponent.}
Suppose the entries of a matrix $A$ commute with one another. Enumerate
the distinct entries as $a_1,\ldots,a_k$, and suppose $a_j^{e_j}=0$.
Then
\begin{equation}\label{a467:pigeon}
 A^r=0,\qquad r=1+\sum_{j=1}^k(e_j-1).
\end{equation}

\emph{Proof.} A product of $r$ entries has multiplicities $m_j$ with
$\sum m_j=r$. Some $m_j\ge e_j$, by the pigeonhole principle. Commuting
the factors groups together $e_j$ copies of $a_j$, so this product is
zero. Now use the entry formula from Lemma 467.1.\hfill$\square$

This proves the matrix assertion for commutative nil rings, even when the
ring is not nilpotent. It also explains exactly why applying the ordinary
commutative multinomial argument to a noncommutative matrix power fails:
the offending repeated letters need not occur consecutively, and moving
them past intervening letters may change the product.

A related triangular subcase needs no commuting assumption. For
$T=\left(\begin{smallmatrix}a&b\\0&d\end{smallmatrix}\right)$,
direct multiplication gives
\[
 (T^n)_{12}=\sum_{j=0}^{n-1}a^jbd^{n-1-j}.
\]
If $a^r=d^s=0$, then $T^{r+s}=0$: in every displayed summand either
$j\ge r$ or $n-1-j\ge s$. This argument works because the triangular
shape forces every word to have only one switch between its two diagonal
states. A general matrix permits repeated switches.

\paragraph{A finite-dimensional nil algebra is nilpotent.}
Here is a proof of a second complete special case without conflating
individual exponents with a ring exponent. Let $A$ be a finite-dimensional
nil algebra over a field $F$. Put $U=F1+A$. The spaces
$A\supseteq A^2\supseteq\cdots$ form a descending chain and therefore
stabilize. Suppose $I=A^r=A^{r+1}$ at some stage.

The left $U$-module $I$ is finitely generated, since an $F$-basis generates
it. If it were nonzero, choose a generating list $v_1,\ldots,v_k$ of
minimal positive length. The equality $I=AI$ gives
\[
 v_k=\sum_{j=1}^k a_jv_j,\qquad a_j\in A.
\]
To obtain these coefficients, first express $v_k$ as a finite sum of
products of elements of $A$ and $I$, and then expand each element of $I$
in the chosen generators; $A$ is a two-sided ideal of $U$.
Because $a_k$ is nilpotent, $1-a_k$ has inverse
$1+a_k+\cdots+a_k^{e-1}$ for some $e$. Thus $v_k$ is generated by
$v_1,\ldots,v_{k-1}$, contradicting minimality. Hence $I=0$.

We conclude that $A$ is nilpotent. In fact its powers must strictly
decrease in dimension until reaching zero, so $A^{\dim_F A+1}=0$.
This proves the matrix statement for finite-dimensional nil algebras by
Lemma 467.1. An infinite-dimensional finitely generated nil algebra does
not satisfy this descending-chain argument. No finite-dimensional
approximation is substituted for it.

\paragraph{A concrete nil ring with no common exponent.}
Fix a field $F$ and form the algebraic direct sum
\[
 R=\bigoplus_{m\ge1}tF[t]/(t^{m+1}).
\]
An element has finite support. A common exponent for the finitely many
nonzero components annihilates it, so $R$ is nil. But in the $m$th
component the class of $t$ has nonzero $m$th power. Thus $R$ is not
nilpotent and has no uniform element exponent either.

Nevertheless every finite set of elements has support in finitely many
components, whose direct sum is a nilpotent ring. Hence $R$ is locally
nilpotent, and every matrix over it is nilpotent. This is a useful stress
test for the conclusion: the conjecture asks for individual matrix
nilpotence, which does not require a common exponent for all matrices.
Trying larger components in this example cannot produce a counterexample.

Contrast the inverse system of finite nilpotent rings
$tF[t]/(t^{m+1})$. Their compatible sequence represented by $t$ is not
nilpotent in $tF[[t]]$: every finite power remains nonzero. Thus checking
nilpotence in every finite quotient with exponents tending to infinity
does not prove nilpotence in a limiting ring. Finite tests need a
uniform exponent or a separate infinite argument before that passage.

\paragraph{Lemma 467.3: why a nonnil matrix gives the ideal obstruction.}
Let $N$ be any nil algebra over $F$, and let $U=F1+N$. In $M_2(U)$
consider the two sets of matrices supported in the first and second
columns respectively, with entries in $N$:
\[
 L_1=\left\{\begin{pmatrix}x&0\\y&0\end{pmatrix}:x,y\in N\right\},
 \qquad
 L_2=\left\{\begin{pmatrix}0&x\\0&y\end{pmatrix}:x,y\in N\right\}.
\]
Both are nil left ideals, and $L_1+L_2=M_2(N)$.

\emph{Proof.} Left multiplication by an arbitrary matrix over $U$
preserves the specified zero column and keeps the other entries in the
two-sided ideal $N$. Additive closure is immediate. For the first shape,
\[
 \begin{pmatrix}x&0\\y&0\end{pmatrix}^{m}
 =\begin{pmatrix}x^m&0\\yx^{m-1}&0\end{pmatrix}.
\]
If $x^e=0$, power $e+1$ is zero. For the second shape, the analogous
formula is
$\left(\begin{smallmatrix}0&xy^{m-1}\\0&y^m\end{smallmatrix}\right)$,
which vanishes at power $e+1$ when $y^e=0$. Finally every matrix over
$N$ is the sum of its two columns.\hfill$\square$

This establishes directly the direction of the classical equivalence
needed to interpret any proposed nonnil matrix. Taking opposite rings
converts left ideals to right ideals, so the side convention does not
remove the obstruction. The construction does not show that the two
column ideals are nilpotent ideals: their individual element exponents
can be unbounded.

\paragraph{Why nilpotent entries do not suffice.}
There is a small exact example showing why the hypothesis must concern
the entire ring generated by the entries. Over any field let
$a=E_{12}$ and $b=E_{21}$ in $M_2(F)$. Both satisfy $a^2=b^2=0$, but
$ab=E_{11}$ and $ba=E_{22}$ are nonzero idempotents. The block matrix
\[
 B=\begin{pmatrix}0&a\\b&0\end{pmatrix}
 \quad\text{satisfies}\quad
 B^{2m}=\begin{pmatrix}E_{11}&0\\0&E_{22}\end{pmatrix}\ne0
 \quad(m\ge1).
\]
It is not nilpotent despite every block entry being nilpotent. This is
not a counterexample to K\"othe: the ambient entry ring is not nil,
and $a+b$ already squares to the identity. The example isolates the
invalid inference from individual entry powers to mixed-word powers.

For the same reason, a computer search over matrices whose entries each
pass a nilpotence test cannot certify that their generated algebra is nil.
One must check every algebra element, an infinite requirement unless a
separate structural theorem reduces it to finite data. The finite-dimensional
lemma shows that a finite-dimensional algebra over a field will not
produce the desired nil-but-matrix-nonnil phenomenon.

\paragraph{A precise infinite-word obstruction remains.}
For a matrix $A$, its $(i,j)$ entry at power $n$ sums labels of all
length-$n$ walks from $i$ to $j$ in a finite directed graph whose edge
labels are its entries. Local nilpotence kills every sufficiently long
word and therefore every such sum. Nilness only kills powers of each
fixed algebra element; it does not directly kill the walk sums uniformly
as the word length grows.

Even nonvanishing of an individual walk word would not alone certify a
nonzero matrix entry: different words can cancel in the algebra. A valid
nonnil witness needs, for arbitrarily large powers, a functional, normal
form, representation, or other invariant proving that the full relevant
sum survives. Conversely a general proof of nilness would need a uniform
annihilation argument for all entries of a fixed matrix. Neither step is
supplied by truncating words at a finite length.

This is the boundary of the present attack. The new source [E1] claims
to overcome it by an explicit infinite algebra construction; this document
has verified the theorem's scope and the elementary ideal implication,
but has not independently certified the complete construction and its
all-elements nilness proof. That distinction is essential when reporting
the current status after a newly posted manuscript.

\paragraph{Finite checks and outcome.}
The companion exact-arithmetic script checks the triangular power formula,
the two column-power formulas, and the block-matrix stress test over a
small finite field. It also enumerates all matrices over a small truncated
polynomial nil algebra and verifies the predicted nilpotence bound.
These tests diagnose the displayed calculations; none purports to verify
nilness of the infinite algebra in [E1].

\textbf{Outcome of this attempt: UNRESOLVED, with an external counterexample
claim explicitly recorded.} Complete positive subcases and exact reductions
are proved above. A primary September 2026 manuscript now claims a negative
answer to the exact catalogue conjecture. The outstanding task for a
definitive independent assessment of that claim is a full audit of its
infinite construction, not further finite-dimensional enumeration or
an assertion that the conjecture's historical open status is unchanged.


\subsection{Sources}
\begin{itemize}
\item[S1]Agata Smoktunowicz, "Some results in noncommutative ring theory," Proceedings of the ICM, Madrid 2006, Vol. II, EMS, pp. 259-269.
\url{https://www.mathunion.org/fileadmin/ICM/Proceedings/ICM2006.2/ICM2006.2.ocr.pdf}.
\textit{Formulation source.}
\item[E1]Tom Adamczewski, Bernhard B\"ohmler, and Ren\'e Marczinzik,
\emph{A counterexample to K\"othe's conjecture and a question of Rowen},
arXiv:2609.07996v1, submitted 7 September 2026, Theorem 1.1.
\url{https://arxiv.org/html/2609.07996v1};
\url{https://arxiv.org/abs/2609.07996}.
\textit{Primary theorem statement and version history inspected
27 September 2026. This is an external counterexample claim; its full
proof is not claimed independently verified here.}
\item[E2]Peter K\'alnai and Jan \v{Z}emli\v{c}ka,
\emph{Self-injective von Neumann regular rings and K\"othe's Conjecture},
arXiv:1912.12159 (2019).
\url{https://arxiv.org/abs/1912.12159}.
\textit{Primary abstract checked for the classical formulation and
countability context; historical status predates [E1].}

\end{itemize}
\end{document}
